% =========================
% Chapter 3
% =========================
\chapter{Related Work}
\label{ch:related_work}

\section{Assistive Vision Systems for Blind Users}
\label{sec:assistive_vision_systems}

Assistive vision systems for blind and visually impaired users have been an active
area of research for several decades, driven by advances in computer vision, wearable
sensing, and audio-based interaction. These systems typically employ cameras mounted
on wearable or handheld devices to capture visual information from the user's
environment and translate it into non-visual feedback, most commonly through speech
or spatialized audio. The primary objective is to enhance environmental awareness,
enabling users to perceive obstacles, objects, and spatial layout that would
otherwise be inaccessible through traditional mobility aids alone.

Early assistive vision approaches focused on specific perceptual functions such as
obstacle detection, depth estimation, and object recognition, often providing simple
auditory cues or alerts to indicate the presence of
hazards~\cite{dakopoulos2010wearable}. As computer vision techniques matured, more
comprehensive systems emerged that attempted to describe broader aspects of the
surrounding scene, including object categories, relative positions, and semantic
context~\cite{li2018vision}. These systems generally rely on continuous or
near-continuous visual processing pipelines, transforming camera input into spoken
descriptions intended to substitute or augment visual perception.

More recent work has explored increasingly sophisticated capabilities, including
semantic scene understanding, indoor layout recognition, and landmark
identification~\cite{fernandes2019review}. Advances in deep learning have enabled
higher recognition accuracy and robustness in complex indoor environments, while
improvements in wearable hardware have facilitated hands-free operation through smart
glasses and body-mounted cameras. Many of these systems emphasize real-time feedback
and aim to provide a rich representation of the user's surroundings, often
prioritizing completeness of information over selectivity.

Despite significant progress, existing assistive vision systems vary widely in how
perception is activated, how information is structured, and how feedback is delivered
to users. This diversity reflects differing assumptions about user needs and system
goals, and it motivates a closer examination of interaction paradigms, information
delivery strategies, and control structures, discussed in the following sections.

\section{Interaction Paradigms: Continuous Narration vs.\ On-Demand Assistance}
\label{sec:interaction_paradigms}

A central design dimension in assistive vision systems is the interaction paradigm
that governs when perception is activated and how information is delivered. Prior
work largely falls into two dominant paradigms: continuous scene narration and
on-demand, query-based assistance.

Continuous narration systems aim to provide persistent situational awareness by
continuously processing visual input and verbalizing detected objects, obstacles, and
scene elements~\cite{li2018vision, fernandes2019review}. While this approach can
increase environmental coverage, it often results in verbose and redundant feedback,
requiring users to continuously filter relevant information from irrelevant details.
Studies in assistive technology and auditory interfaces have shown that such
uninterrupted information streams can significantly increase cognitive load,
interfere with attention to ambient sounds, and reduce overall task
performance~\cite{martinez2014cognitive, hossain2010cognitive}.

In contrast, on-demand and query-based assistive systems restrict perception and
feedback to moments explicitly initiated by the user through individual
queries~\cite{khan2020insight, li2018vision}. By limiting feedback to user-triggered
queries, on-demand approaches reduce continuous cognitive burden and offer greater
control over information flow. However, they introduce a different challenge: users
must know what to ask for and when. In unfamiliar environments, blind users may lack
sufficient situational context to formulate appropriate queries, leading to delayed
or incomplete assistance.

Both paradigms reflect implicit assumptions about user behavior and system
responsibility. Continuous narration assumes that more information leads to better
awareness; query-based systems assume that users can effectively manage their own
information needs. In practice, neither assumption consistently holds. This tension
highlights a gap in current assistive vision research: the lack of interaction models
that dynamically align perception and feedback with user intent and task context.
Lumen's task-scoped paradigm addresses exactly this gap --- perception is activated
by a task, structured by that task's phases, and terminated with it --- and its
navigation mode goes one step further by treating the \emph{user} as a collaborative
resource, asked for a direction hint only when perception alone is genuinely
ambiguous.

\section{Indoor Navigation and Object-Finding Systems}
\label{sec:indoor_navigation_object_finding}

Indoor navigation and object-finding represent two of the most common and challenging
assistive tasks. Unlike outdoor navigation, which can rely on GPS and large-scale
mapping, indoor environments are highly structured, dynamic, and GPS-denied,
requiring alternative sensing and localization strategies~\cite{fernandes2019review}.

Indoor navigation systems for blind users often guide users between predefined
locations using landmarks, visual markers, or spatial cues extracted from the
environment~\cite{guerrero2012navigation, rodriguez2012navigation}. Notably,
infrastructure-based systems such as NavCog~\cite{ahmetovic2016navcog} achieve
reliable turn-by-turn guidance by instrumenting the building with beacons and
pre-built maps --- an approach that works well in flagship venues but does not scale
to arbitrary homes and offices. Vision-based approaches avoid instrumentation by
detecting doors, corridors, and signage from camera input, but many rely on prior
mapping or fiducial markers, and most assume a known route rather than genuine
exploration of an unknown space. Lumen's navigation deliberately occupies the
uninstrumented end of this spectrum: no maps, no markers, no SLAM --- exploration
driven by door detection and semantic room recognition, with the human contributing
coarse direction knowledge that the sensors cannot provide.

Object-finding systems help users locate specific items within their immediate
surroundings~\cite{brilhault2011object, tapu2014wearable}, commonly employing
real-time object detection and audio cues to guide users toward a target. Although
object-finding is often treated as a simpler problem than navigation, it presents
similar challenges in perception reliability, spatial feedback, and interaction flow.
Many existing systems continuously scan the environment or require repeated user
queries, offering limited structure for task initiation, progression, and
termination. Lumen structures the entire object-finding interaction as one task with
explicit phases --- search, track, reach, grasp --- ending either by the user's
confirmation or by the system detecting the grasp itself.

\section{Cognitive Load and Information Management}
\label{sec:cognitive_load}

Cognitive load has been widely recognized as a critical factor in the effectiveness
of assistive technologies. Assistive systems rely on auditory feedback as the primary
communication channel, requiring users to process spoken information while
simultaneously attending to environmental sounds, spatial cues, and physical
navigation. Excessive or poorly structured auditory information overwhelms users,
reducing task performance and situational
awareness~\cite{hossain2010cognitive, martinez2014cognitive}.

Beyond the quantity of information, prior work emphasizes timing, structure, and
context: information presented outside an immediate task context is less likely to be
retained and more likely to interfere with concurrent cognitive
processes~\cite{wickens2008multiple}. Despite these insights, many assistive vision
systems treat cognitive load as a secondary concern, addressed through ad hoc
verbosity reduction rather than architectural mechanisms that regulate when
perception is active and when information is delivered.

Lumen treats cognitive load management as a first-class design objective at three
levels: the task-scoped paradigm bounds \emph{when} the system may speak at all; a
two-tier speech policy distinguishes priority announcements from droppable ambient
nudges, so guidance never queues up and lags behind reality; and utterance merging
ensures that no spoken line is ever interrupted by the next one.

\section{Wearable Platforms and Voice-Based Interaction}
\label{sec:wearable_voice}

Wearable platforms --- smart glasses, head- or chest-mounted cameras, belt-worn
systems --- have been the dominant form factor for assistive vision, as they enable
hands-free sensing~\cite{tapu2014wearable, dakopoulos2010wearable}. Their design is
constrained by size, weight, battery life, and computational resources; consequently,
many systems offload processing to paired smartphones or cloud
services~\cite{chen2015wearable}, trading local constraints for latency and
connectivity challenges.

Voice-based interaction has emerged as the natural control modality for such
systems~\cite{porzi2018assistive}, but speech interfaces require careful design to
avoid ambiguity, misrecognition, and dialogue overhead, particularly in noisy indoor
environments. Moreover, many systems treat voice as a thin control layer mapping
commands directly to isolated functions, with limited consideration of interaction
flow, task progression, or system state.

Lumen draws a pragmatic conclusion from this literature: the modern smartphone
\emph{already is} the wearable --- camera, microphone, speaker, IMU, compass, and
connectivity in one ubiquitous device --- and delivering the client as a web page
removes even the app-install barrier. The engineering effort thus shifts from
hardware to interaction structure: an explicit state machine governs how tasks are
initiated, executed, interrupted, and terminated by voice.

\section{System Architectures, Control Models, and Safety}
\label{sec:architectures_control_safety}

The effectiveness of assistive technologies depends not only on sensing accuracy and
interaction modalities, but also on the underlying architecture that governs task
execution, control flow, and error handling. Prior work commonly describes functional
components --- perception, localization, feedback --- while leaving the coordinating
control logic implicit~\cite{tapu2014wearable, dakopoulos2010wearable}. Many systems
adopt reactive architectures where sensory inputs directly trigger
responses~\cite{porzi2018assistive}; without explicit state, perception may remain
active indefinitely, feedback may repeat unpredictably, and behavior under edge cases
is hard to reason about.

Safety and failure handling are particularly critical for blind users, for whom
incorrect guidance has immediate physical consequences, yet most systems rely on ad
hoc safeguards such as confidence thresholds~\cite{fernandes2019review}. In robotics
and human--computer interaction, state-based control models --- finite state machines
and their descendants --- have long provided deterministic behavior, clear task
boundaries, and predictable recovery~\cite{brooks1986robust,
colledanchise2018behavior}, but have seen limited adoption in assistive vision
systems.

Lumen adopts explicit state-machine control at both the session level and inside the
navigation task, and extends the safety-first philosophy into perception itself:
every navigation-critical percept must survive multiple independent checks (temporal
consistency, geometric gates, semantic verification) before it is ever spoken to the
user, and obstacle warnings preempt all other guidance while the user is walking.
Because the decision logic is pure and clock-injected, these safety behaviors are not
just design intentions --- they are verified by automated tests.

\section{Summary and Research Gap}

Across this literature, three gaps recur. First, interaction paradigms are polarized
between continuous narration and reactive querying, with little work on structuring
assistance around explicit, bounded tasks. Second, uninstrumented indoor navigation
--- no maps, no beacons, no SLAM --- remains largely unaddressed for blind users,
despite being the common case in homes and unfamiliar buildings. Third, control
logic and safety handling are typically implicit, making system behavior unpredictable
and unverifiable. Lumen's design responds to all three: task-scoped perception under
an explicit FSM, goal-directed exploration navigation that fuses semantic room
recognition with door detection and human hints, and a verification-first
architecture whose behavioral guarantees are enforced by an automated test suite.
